% Generated draft: requires mathematical and source review.
\documentclass[11pt]{article}
\usepackage[margin=25mm]{geometry}
\usepackage{fontspec}
\setmainfont{Latin Modern Roman}
\usepackage{amsmath,amssymb,mathtools,unicode-math}
\setmathfont{Latin Modern Math}
\usepackage{enumitem,needspace,xurl,hyperref}
\setlength{\parindent}{0pt}
\setlength{\parskip}{5pt}
\setlength{\emergencystretch}{3em}
\newcommand{\N}{\mathbb{N}}
\newcommand{\Z}{\mathbb{Z}}
\newcommand{\Q}{\mathbb{Q}}
\newcommand{\R}{\mathbb{R}}
\newcommand{\C}{\mathbb{C}}
\newcommand{\F}{\mathbb{F}}
\newcommand{\A}{\mathbb{A}}
\newcommand{\PP}{\mathbb{P}}
\newcommand{\E}{\mathbb{E}}
\newcommand{\eps}{\varepsilon}
\newcommand{\dd}{\,\mathrm{d}}
\begin{document}
\section*{UnsolvedMath ID: 801119}
\textbf{Source problem code:} catalogue entry
\subsection{Definitions and mathematical statement}
A type \(\mathrm{II}_1\) factor is an infinite-dimensional von Neumann algebra \(M\) with center \(\C 1\) and a faithful normal tracial state \(\tau_M\), normalized by \(\tau_M(1)=1\). Write \(\|x\|_{2,M}=\tau_M(x^*x)^{1/2}\). In the finite-set formulation, \(M\) has property Gamma if, for every finite subset \(F\subseteq M\) and every \(\varepsilon>0\), there is a unitary \(u\in M\) such that
\[
\tau_M(u)=0,\qquad \|ux-xu\|_{2,M}<\varepsilon\quad(x\in F).
\]
Let \(R\) denote the separable hyperfinite type \(\mathrm{II}_1\) factor. A factor \(M\) is McDuff if \(M\cong M\,\overline\otimes\,R\), where \(\overline\otimes\) is the von Neumann algebra tensor product. In the standard continuous first-order language of tracial von Neumann algebras, write \(M\equiv N\) when every sentence has the same value in \(M\) and \(N\).

The source title supplies the type \(\mathrm{II}_1\) context, but the supplied excerpt does not specify whether the factors must have separable predual. Let \(\mathscr K\) denote the intended class of type \(\mathrm{II}_1\) factors, with this separability restriction left unresolved, and set
\[
\mathscr C_{\mathscr K}=\{M\in\mathscr K: M\text{ has property Gamma and is not McDuff}\}.
\]
The two questions are:
\begin{enumerate}
\item Do there exist \(M,N\in\mathscr C_{\mathscr K}\) such that \(M\not\equiv N\)?
\item Does there exist a family \((M_\alpha)_{\alpha<2^{\aleph_0}}\) in \(\mathscr C_{\mathscr K}\) such that
\[
\alpha\ne\beta\ \Longrightarrow\ M_\alpha\not\equiv M_\beta?
\]
\end{enumerate}
The second question interprets ``continuum many such examples'' as continuum many pairwise elementarily inequivalent factors. This interpretation and the intended separability convention require confirmation from the source. If nonseparable factors are intended, the finite-set definition of property Gamma above should also be checked against the source's convention.

The supplied statement carries a solved annotation attributing a solution to J. Peterson, whereas the registry status is \texttt{open}. These are conflicting supplied status labels; no resolution or mathematical verification is asserted here.
\subsection{Short English statement}
The problem asks whether two type \(\mathrm{II}_1\) factors can both have property Gamma and fail to be McDuff while having different elementary theories, and whether continuum many pairwise different theories occur among such factors. The separability convention is unspecified. The supplied solved annotation conflicts with the registry's open status.
\subsection{Sources}
\begin{itemize}
\item Local UnsolvedMath registry, supplied complete record with permanent ID \(801119\), source problem code \texttt{catalogue entry}, titled ``Ultraproduct \(\mathrm{II}_1\) Factors --- Problem 20.'' The registry attributes the problem to Elayavalli, with proposed year 2025.
\item Supplied source attribution: Elayavalli, \emph{50 Open Problems: Ultraproduct \(\mathrm{II}_1\) Factors} (2025), Problem 20, \url{https://arxiv.org/abs/2511.20377}. The record reports access on 2026-10-16 and extraction by \texttt{arxiv\_latex}; these are supplied provenance claims, not independently verified access or source verification.
\item The supplied statement includes the annotation ``Solved by J. Peterson'' with reference label \texttt{peterson2026ultraproductapproximationspropertyt}. No full bibliographic entry or URL for that reference is supplied, so the attribution and the scope of the claimed solution remain unverified.
\end{itemize}
\end{document}
